\section{Synthetic study: where the adjudication goes}
\label{sec:synthetic}

\subsection{Definitions}
\label{sec:metrics}

The confusions in this paper are defined over the \emph{verdict}, not over the
MR's trigger, because the trigger is evidence and the verdict is what a
developer is told.

\begin{center}
\begin{tabular}{lll}
\toprule
ground truth & verdict & counted as \\
\midrule
violation & \texttt{violation-proven} & true positive \\
violation & not \texttt{violation-proven} (including abstention) & false negative \\
correct & \texttt{violation-proven} & \textbf{false proof} (FP) \\
correct & not \texttt{violation-proven} & true negative \\
\bottomrule
\end{tabular}
\end{center}

A \emph{false proof} is therefore the strict subset of false alarms that have
been promoted to \texttt{violation-proven}: an error the tool asserts as
established. We also report two rates that this table does not capture. The
\emph{detection-keep rate} is the share of true violations on which the MR
fired and the method still alarmed; its denominator is restricted to
true violations on which the MR fired, so that it measures whether a correction
\emph{discarded} a success rather than whether detection was possible. The
\emph{abstention count} counts only abstentions justified by an unreadable
oracle, not cases where the MR simply did not fire.

\subsection{The switch that cannot be thrown both ways}

Table~\ref{tab:grid} crosses the visibility of the sharing channel against the
correctness of the application's authorization, on the representative slice.

\begin{table}[t]
\centering
\caption{The 2$\times$2 grid on the representative slice. FP~=~false proofs,
FN~=~false negatives; each cell has \ResBEThreevulnZerounlistedCases{} records
over \ResDesignCells{} design cells. Intervals are Wilson 95\% and, per
Section~\ref{sec:stats}, are conventional rather than sampling-based.}
\label{tab:grid}
\begin{tabular}{llrrr}
\toprule
channel & application & FP & FN & FP rate [95\% CI] \\
\midrule
unlisted & correct   & \ResBEThreevulnZerounlistedFp{} & \ResBEThreevulnZerounlistedFn{}
         & \ResBEThreevulnZerounlistedFpPct{}\% [\ResBEThreevulnZerounlistedFpLo{}, \ResBEThreevulnZerounlistedFpHi{}] \\
listed   & correct   & \ResBEThreevulnZerolistedFp{} & \ResBEThreevulnZerolistedFn{}
         & \ResBEThreevulnZerolistedFpPct{}\% [\ResBEThreevulnZerolistedFpLo{}, \ResBEThreevulnZerolistedFpHi{}] \\
unlisted & violating & \ResBEThreevulnOneunlistedFp{} & \ResBEThreevulnOneunlistedFn{}
         & \ResBEThreevulnOneunlistedFpPct{}\% [\ResBEThreevulnOneunlistedFpLo{}, \ResBEThreevulnOneunlistedFpHi{}] \\
listed   & violating & \ResBEThreevulnOnelistedFp{} & \ResBEThreevulnOnelistedFn{}
         & \ResBEThreevulnOnelistedFpPct{}\% [\ResBEThreevulnOnelistedFpLo{}, \ResBEThreevulnOnelistedFpHi{}] \\
\bottomrule
\end{tabular}
\end{table}

Three of the four cells contain an error
(\ResGridCellsWithError{}/\ResGridCellsTotal{}). One cell is clean---
listed$\times$correct, with no error at all---but it is clean for a trivial
reason, and we say so: on that cell the MR never fires, so there is nothing for
the adjudicator to get wrong, and there is no true violation to miss either. The
same switch that makes it clean makes the paired
listed$\times$violating cell err in the other direction, where the MR fails to
fire on a genuine violation and produces \ResBEThreevulnOnelistedFn{} false
negatives out of \ResBEThreevulnOnelistedTp{}+\ResBEThreevulnOnelistedFn{}.

The substantive reading of Table~\ref{tab:grid} is therefore not ``cells are
error-free'' but: \emph{avoiding false proofs and keeping detection are
controlled by one switch, and the switch cannot be thrown both ways at once}.
On unlisted cells the false-proof rate is
\ResBEThreevulnZerounlistedFpPct{}\% with no false negatives; on listed cells the
false-proof rate is \ResBEThreevulnZerolistedFpPct{}\% with false negatives
appearing. Section~\ref{sec:synthetic-dilemma} shows the same tension at the
level of individual scenarios.

\subsection{Attribution is clean, and the anchor is small}

On the anchored configuration the false-proof rate is
\ResFpEthreeHits{}/\ResFpEthreeTotal{} = \ResFpEthreePct{}\%, and this paper will
never quote it without its interval, which is
$[\ResFpEthreeLo{}\%, \ResFpEthreeHi{}\%]$. The two controls that isolate
attribution are clean: with the legitimate-sharing exception absent the rate is
\ResBaseFpRate{} (\ResBaseFpHits{}/\ResBaseFpTotal{}), and with the channel
listed the rate is \ResListedCtlFpRate{}
(\ResListedCtlFpHits{}/\ResListedCtlFpTotal{}). The phenomenon is therefore
attributable to the presence of a legitimate sharing exception that the MR
cannot know about, not to the MR merely being executed. Recall on the unlisted
violation scenarios S2, S4 and S5 is \ResRecallSTwo{},
\ResRecallSFour{} and \ResRecallSFive{} respectively.

Because \ResFpEthreeTotal{} is two cases, we lead elsewhere with the larger
counts: the per-cell \ResBEThreevulnZerounlistedFpPct{}\% of
Table~\ref{tab:grid} and, in Section~\ref{sec:synthetic-ablation}, a
\ResSynRecords{}-record panel.

\subsection{Filtering exceptions does not help}
\label{sec:synthetic-ladder}

One natural repair is to add exception predicates: recognise legitimate sharing
and filter it. We evaluated the exception ladder---raw, then one exception
level, then two---on the anchored false-proof rate. The rate is
\ResLadderrawRate{}, \ResLaddereOneRate{} and \ResLaddereOneeTwoRate{}: it does
not move. This is a constructive result rather than a disappointing one: the
exception levels we can express are themselves functions of the same
reachability model, so they inherit its blind spot. Adding predicates from the
same family cannot repair a failure caused by that family.

\subsection{The framework cannot see the channel that matters}
\label{sec:synthetic-observability}

Table~\ref{tab:obs} separates the two channels. When the legitimate share was
created through the framework's own facility, the framework's log holds
\ResObsFwLogRecords{} records on the scenario, and the exception-aware
adjudicator v2 produces \ResObsFwVtwoFalseProof{} false proofs
(rate \ResObsFwVtwoFpRate{}). When the same legitimate share was created
\emph{out of band}, the framework's log holds \ResObsOobLogRecords{} records,
and v2's false-proof count rises to \ResObsOobVtwoFalseProof{} of
\ResObsOobFalseAlarm{} alarms, i.e.\ a rate of \ResObsOobVtwoFpRate{}, with
every one of them promoted to \texttt{violation-proven}.

\begin{table}[t]
\centering
\caption{Observability of the sharing channel. Each scenario has
\ResObsOobFalseAlarm{} alarms over \ResDesignCells{} design cells. A log record
count of \ResObsOobLogRecords{} means the framework has no evidence that the
share exists.}
\label{tab:obs}
\begin{tabular}{lrrrr}
\toprule
scenario & framework log records & alarms & v2 false proofs & v2 rate \\
\midrule
framework-created legitimate share (S3) & \ResObsFwLogRecords{} & \ResObsFwFalseAlarm{} & \ResObsFwVtwoFalseProof{} & \ResObsFwVtwoFpRate{} \\
out-of-band legitimate share (S6)       & \ResObsOobLogRecords{} & \ResObsOobFalseAlarm{} & \ResObsOobVtwoFalseProof{} & \ResObsOobVtwoFpRate{} \\
out-of-band share with a real violation (S7) & \ResObsOobBolaLogRecords{} & \ResObsOobBolaFalseAlarm{} & \ResObsOobBolaVtwoFalseProof{} & \ResObsOobBolaVtwoFpRate{} \\
\bottomrule
\end{tabular}
\end{table}

This is boundary 1 in its sharpest form.

\begin{figure}[t]
\centering
\includegraphics[width=0.98\linewidth]{fig2_scenarios.pdf}
\caption{Errors are scenario-structured. The bookkeeping adjudicator's false
positives concentrate on the out-of-band legitimate-share scenario (S6), where
the framework has no log record of the share, and the state-calibrated
adjudicator's false negatives concentrate where its reachability precondition
blocks a detection. Every one of the \ResAbDisagreementN{} records on which
\texttt{direct\_observe} and v3 disagree is of the second kind.}
\label{fig:scenarios}
\end{figure}

The MR cannot distinguish ``this
share does not exist'' from ``this share exists but was not created through the
interface I crawled''. Its only available representation of the first statement
is the second, so an out-of-band grant is reported as a proven violation. Note
also that on S7---where there \emph{is} a genuine violation behind the
out-of-band share---the same blindness would make the alarm right for the wrong
reason: of the \ResObsOobBolaBolaTotal{} true violations on that scenario, the
exception-aware method reaches \ResObsOobBolaVtwoProvenBola{} and the
state-calibrated method \ResObsOobBolaVthreeProvenBola{}. Being right by
accident is not detection, and we do not count it as such. Note also that on
S7 the alarms and the true violations coincide numerically, which is exactly
what makes the accident invisible to an input-level metric.

\subsection{One switch, two scenarios}
\label{sec:synthetic-dilemma}

The grid's tension can be exhibited inside a single MR, with no change other
than the channel-visibility switch. Consider S3, where the correct behaviour is
\emph{not} to alarm (there is a legitimate share and no violation), and S5,
where the correct behaviour \emph{is} to alarm (a read-only share must not
authorise a write). Under the unlisted configuration the MR is correct on S5
but wrong on S3; under the listed configuration it is correct on S3 but wrong on
S5. Writing $c_x(S)$ for ``configuration $x$ is correct on scenario $S$'': we
obtain $c_{\text{unlisted}}(\text{S3}) = \ResSthreeSfiveRowZeroSthreeOk{}$,
$c_{\text{unlisted}}(\text{S5}) = \ResSthreeSfiveRowZeroSfiveOk{}$,
$c_{\text{listed}}(\text{S3}) = \ResSthreeSfiveRowOneSthreeOk{}$ and
$c_{\text{listed}}(\text{S5}) = \ResSthreeSfiveRowOneSfiveOk{}$, so no
configuration in this family is correct on both---the disjunction over
configurations is \ResSthreeSfiveAnyCorrect{}. For this MR the choice of switch
is a choice of \emph{which} error to make, not whether to make one.

\subsection{The aggregation hides the regime change}

Table~\ref{tab:aggregate} aggregates verdict quality over all configurations,
using the recovery rate (share of raw alarms the method resolves correctly) and
the detection-keep rate.

\begin{table}[t]
\centering
\caption{Verdict quality over all configurations. ``Recovery'' is the share of
raw alarms resolved correctly; ``keep'' is the share of fired true violations
that still alarm. \texttt{degrade\_all} is the floor that abstains on
everything.}
\label{tab:aggregate}
\begin{tabular}{lrrr}
\toprule
method & false alarms & recovery & keep \\
\midrule
v1 (raw)                 & \ResVdVOneFalseAlarm{}       & \ResVdVOneRecoveryRate{}       & \ResVdVOneKeepRate{} \\
v2 (exception-aware)     & \ResVdVTwoFalseAlarm{}       & \ResVdVTwoRecoveryRate{}       & \ResVdVTwoKeepRate{} \\
v3 (state-calibrated)    & \ResVdVThreeFalseAlarm{}     & \ResVdVThreeRecoveryRate{}     & \ResVdVThreeKeepRate{} \\
\texttt{degrade\_all}    & \ResVdDegradeallFalseAlarm{} & \ResVdDegradeallRecoveryRate{} & \ResVdDegradeallKeepRate{} \\
\bottomrule
\end{tabular}
\end{table}

Read naively, v3 looks like the best method in the table: it recovers
\ResVdVThreeRecoveryRate{} of raw alarms while \texttt{degrade\_all} keeps
\ResVdDegradeallKeepRate{}. That reading is wrong, and the way to see that it is
wrong is to stratify. v3's behaviour changes regime depending on whether an
authorization introspection surface is readable, and the aggregate averages two
different regimes:

\begin{center}
\begin{tabular}{lrrrr}
\toprule
stratum & records & false proofs & keep & abstentions \\
\midrule
introspection readable & \ResIntroOnN{}  & \ResIntroOnVthreeFpN{} & \ResIntroOnVthreeKeepRate{} & \ResIntroOnAbstain{} \\
introspection absent   & \ResIntroOffN{} & \ResIntroOffVthreeFpN{} & \ResIntroOffVthreeKeepRate{} & \ResIntroOffAbstain{} \\
\bottomrule
\end{tabular}
\end{center}

On the readable stratum v3 is genuinely good: it removes the false proofs the
exception-aware method still produces (\ResIntroOnVthreeFpN{} against v2's
\ResIntroOnVtwoFpN{}) and still alarms on all \ResIntroOnVthreeProvenBola{}
fired true violations. On the absent stratum its keep rate is
\ResIntroOffVthreeKeepRate{}: it alarms on none of the
\ResIntroOffVtwoProvenBola{} true violations the exception-aware method still
alarms on, and it abstains \ResIntroOffAbstain{} times. The aggregate reports
this as the best row in the table. It is not better; it is a different thing. A
method that abstains on everything is not accurate, and any report that pools
these two strata turns honest abstention into apparent performance. We therefore
state, as a reporting rule, that v3's numbers may appear only stratified.

\begin{figure}[t]
\centering
\includegraphics[width=0.98\linewidth]{fig3_strata.pdf}
\caption{The same method, two regimes. Stratified by readability of the
authorization introspection surface, v3's false proofs vanish only on the
stratum where it can still keep its detections; on the other stratum the zero
false-proof count coincides with a keep rate of
\ResIntroOffVthreeKeepRate{}. Pooling the two strata (Table~\ref{tab:aggregate})
converts abstention into apparent performance.}
\label{fig:strata}
\end{figure}

\subsection{The decisive ablation: does the metamorphic step carry work?}
\label{sec:synthetic-ablation}

Our earlier iteration proposed v3 as a mechanism, and the claim it rested on was
that the metamorphic step was necessary to it. We tested that claim by removing
the metamorphic step and replacing it with the most naive thing that consumes
the same observable quantity.

The distinction matters. A first control, \texttt{direct\_only}, consults the
readable authorization surface but ignores whether the access succeeded; it is a
straw man, and it shows: \ResAbDirectOnlyFP{} false proofs against v3's
\ResAbVThreeFP{}, precision \ResAbDirectOnlyPrec{}. Rejecting \texttt{direct\_only}
proves nothing, because it discards half of the available evidence. The control
that is actually fair is \texttt{direct\_observe}: an ordinary request that
succeeds, together with the same readable authorization surface. It consumes
exactly what the MR-based adjudicator consumes.

\begin{table}[t]
\centering
\caption{End-to-end confusion matrix over all \ResSynRecords{} records
(\ResAbNBola{} true violations, \ResAbNNormal{} correct cases). Intervals are
Wilson 95\% and are conventional (Section~\ref{sec:stats}).
\texttt{direct\_only} is a deliberately unfair control and is shown to make that
point.}
\label{tab:ablation}
\begin{tabular}{lrrrrrrl}
\toprule
method & TP & FP & FN & TN & precision & recall & F1 \\
\midrule
v1 (raw)              & \ResAbVOneTP{} & \ResAbVOneFP{} & \ResAbVOneFN{} & \ResAbVOneTN{} & \ResAbVOnePrec{} & \ResAbVOneRec{} & \ResAbVOneFone{} \\
v2 (exception-aware)  & \ResAbVTwoTP{} & \ResAbVTwoFP{} & \ResAbVTwoFN{} & \ResAbVTwoTN{} & \ResAbVTwoPrec{} & \ResAbVTwoRec{} & \ResAbVTwoFone{} \\
v3 (state-calibrated) & \ResAbVThreeTP{} & \ResAbVThreeFP{} & \ResAbVThreeFN{} & \ResAbVThreeTN{} & \ResAbVThreePrec{} & \ResAbVThreeRec{} & \ResAbVThreeFone{} \\
\texttt{degrade\_all}& \ResAbDegradeAllTP{} & \ResAbDegradeAllFP{} & \ResAbDegradeAllFN{} & \ResAbDegradeAllTN{} & \ResAbDegradeAllPrec{} & \ResAbDegradeAllRec{} & \ResAbDegradeAllFone{} \\
\texttt{direct\_only} & \ResAbDirectOnlyTP{} & \ResAbDirectOnlyFP{} & \ResAbDirectOnlyFN{} & \ResAbDirectOnlyTN{} & \ResAbDirectOnlyPrec{} & \ResAbDirectOnlyRec{} & \ResAbDirectOnlyFone{} \\
\texttt{direct\_observe} & \ResAbDirectObserveTP{} & \ResAbDirectObserveFP{} & \ResAbDirectObserveFN{} & \ResAbDirectObserveTN{} & \ResAbDirectObservePrec{} & \ResAbDirectObserveRec{} & \ResAbDirectObserveFone{} \\
\bottomrule
\end{tabular}
\end{table}

Table~\ref{tab:ablation} is the result that closes this line of work.

\begin{figure}[t]
\centering
\includegraphics[width=0.98\linewidth]{fig4_methods.pdf}
\caption{The MR trigger does not carry the work. Left: the deliberately unfair
control \texttt{direct\_only} compares badly, and rejecting it proves nothing.
Right: the fair control \texttt{direct\_observe}, which consumes the same
observable quantity as the true label, attains the same point-estimate precision
with \ResAbDirectObserveFP{} false proofs and a recall difference of
\ResAbRecallDiff{} (95\% CI
$[\ResAbRecallDiffLo{}, \ResAbRecallDiffHi{}]$, includes zero).}
\label{fig:methods}
\end{figure}

\texttt{direct\_observe} has \ResAbDirectObserveFP{} false proofs against v3's
\ResAbVThreeFP{}, with the same point-estimate precision
(\ResAbDirectObservePrec{}), and its recall is \ResAbDirectObserveRec{} against
v3's \ResAbVThreeRec{}. The recall difference of \ResAbRecallDiff{} is not
distinguishable from zero (95\% CI
$[\ResAbRecallDiffLo{}, \ResAbRecallDiffHi{}]$), so the correct description is
a \emph{descriptive, point-estimate} comparison, not domination; because no
margin was pre-specified, no formal hypothesis-test claim attaches to it, and
the descriptive
comparison is already enough to answer the question, because the two methods
disagree on only
\ResAbDisagreementN{} of the records. On every one of those, v3 does not fire
the MR and \texttt{direct\_observe} alarms on a genuine violation: the
reachability precondition blocked a detection that a plain request finds.

We also report what a selective presentation would have hidden. v3's precision
of \ResAbVThreePrec{} is real, and it was bought with
\ResAbVThreeFN{} false negatives; its F1 is \ResAbVThreeFone{}. Quoting the
precision alone, as our earlier iteration did, would have described a method
that exchanges \ResAbVOneFP{} false alarms for \ResAbVThreeFN{} missed
violations as though it were an improvement. It is not, and it is not a
mechanism that carries work: the metamorphic step adds nothing here that a
single ordinary request does not already provide.

\subsection{Reproducing the reported specificity, and why it cannot see this}

The reported specificity has the form $1 - \text{false alarms} / \text{follow-up
inputs}$. We cannot reproduce the engine's exact input count, so we approximate
the denominator as two follow-up inputs per case; this is enough to show the
denominator's resolution, which is the only point we wish to make.
Table~\ref{tab:spec} gives the result.

\begin{table}[t]
\centering
\caption{The reported specificity convention, approximated. The denominator is
an input count, not a scenario count.}
\label{tab:spec}
\begin{tabular}{lrrr}
\toprule
configuration & cases & false alarms & specificity \\
\midrule
no legitimate-sharing exception & \ResSpecRowZeroCases{} & \ResSpecRowZeroFa{} & \ResSpecRowZeroSpec{} \\
exception, unlisted channel     & \ResSpecRowOneCases{}  & \ResSpecRowOneFa{}  & \ResSpecRowOneSpec{} \\
exception, listed channel       & \ResSpecRowTwoCases{}  & \ResSpecRowTwoFa{}  & \ResSpecRowTwoSpec{} \\
\bottomrule
\end{tabular}
\end{table}

The unlisted configuration---the one where a whole class of legitimate access is
misreported---still yields a specificity of \ResSpecRowOneSpec{}
(\ResSpecRowOneSpecPct{}\%). A denominator that counts generated inputs is
divided by the number of requests the generator emitted, so a failure that is
\emph{uniform across a scenario class} is diluted by every other input in the
run. This is boundary 3: the metric is not wrong, it is
\emph{low-resolution} with respect to the failure mode. It is also not our
finding: the predecessor design attributes its own residual false alarms to
crawler coverage of ``all URLs for all users''~\cite{Mai2019Metamorphic}, and
\emph{that} is the same reachability-blindness we are describing, reported at a
  grain at which it cannot be seen.

\paragraph{A proposition on the metric's resolution}
The dilution above is not incidental; it is structural. Define the
input-level specificity $S_{\mathrm{input}} = 1 - \mathit{FP}_{\mathrm{inputs}} / N_{\mathrm{inputs}}$
and the scenario-level specificity
$S_{\mathrm{scenario}} = 1 - N_{\mathrm{failed}} / N_{\mathrm{scenarios}}$, where
$\mathit{FP}_{\mathrm{inputs}}$ counts false alarms across all generated inputs
and $N_{\mathrm{failed}}$ counts the scenarios in which the failure is uniform.

\noindent\textbf{Proposition.} Fix $f \ge 1$ failed scenarios, each producing a
fixed number $a \ge 1$ of false-alarm inputs, and $s \ge 1$ successful
scenarios each contributing $m$ additional, alarm-free inputs. Then
$S_{\mathrm{input}} = 1 - (f a) / (f a + s m)$, which tends to $1$ as
$m \to \infty$, while $S_{\mathrm{scenario}} = 1 - f / (f + s)$ is independent
of $m$ and strictly less than $1$.

\noindent\textbf{Proof.} The total input count is
$N_{\mathrm{inputs}} = f a + s m$ and the false-alarm count is
$\mathit{FP}_{\mathrm{inputs}} = f a$; substitution gives
$S_{\mathrm{input}} = 1 - (f a) / (f a + s m)$. As $m \to \infty$ the fraction
tends to $0$, so $S_{\mathrm{input}} \to 1$, whereas the scenario-level figure
depends only on the fixed counts $f$ and $s$. \hfill $\square$

\smallskip
\noindent\textbf{Corollary.} No fixed input-level specificity threshold
$t < 1$ can guarantee detection of a scenario-uniform failure class: choosing
$m > (\mathit{fa}/s)\cdot t/(1-t)$ drives $S_{\mathrm{input}}$ above $t$ while
the $f$ failing scenarios remain failing. \hfill $\square$

The metric is therefore not merely imprecise: it is arbitrarily dilutable
with respect to a failure mode that is uniform across a scenario class---adding
unaffected inputs drives it toward $1$ without touching the failure---so it
lacks stable scenario-level resolution.

\subsection{Second implementation}

A second, independently written implementation of the same predicates agreed
with the first on \ResAltAgree{} of \ResAltCompared{} records
(\ResAltAgreeRate{}; Wilson 95\% CI
$[\ResAltAgreeRateLo{}, \ResAltAgreeRateHi{}]$). This bounds
implementation accident only. It does not bound a shared blind spot in the
predicate semantics, because both implementations were written from the same
reading of those semantics, and we say so rather than presenting the agreement
as corroboration.
